% ——————————————————————————————————————————————————————————————————
% Пакеты
% ——————————————————————————————————————————————————————————————————
\usepackage[T2A]{fontenc}
\usepackage[english, main=russian]{babel}
\usepackage{pscyr}
\usepackage{hyperref}
\usepackage{graphicx}
\usepackage{xcolor}
\usepackage{titlesec}
\usepackage{titletoc}
\usepackage{framed}
\usepackage{tikz}
\usepackage[default]{sourceserifpro}
\usepackage[a4paper, margin=2.5cm]{geometry}


% ——————————————————————————————————————————————————————————————————
% Цвета
% ——————————————————————————————————————————————————————————————————
\definecolor{nngblue}{HTML}{1B3A57}
\definecolor{nnggray}{HTML}{555555}
\definecolor{nnglink}{HTML}{385AFF}

\hypersetup{
  colorlinks=true,
  urlcolor=nnglink,
  linkcolor=nnglink,
  citecolor=nnglink
}


% ——————————————————————————————————————————————————————————————————
% Шрифты
% ——————————————————————————————————————————————————————————————————
\renewcommand{\rmdefault}{ftm}
\renewcommand{\sfdefault}{ftx}
\renewcommand{\familydefault}{\rmdefault}

\titleformat{\section}{\Large\bfseries\sffamily}{}{0pt}{}
\titleformat{\subsection}{\large\bfseries\sffamily}{}{0pt}{}


% ——————————————————————————————————————————————————————————————————
% Круглое фото автора
% ——————————————————————————————————————————————————————————————————
\newcommand{\circlephoto}[1]{%
  \begin{tikzpicture}
    \path[use as bounding box] (-20bp,-20bp) rectangle (20bp,20bp);
    \clip (0,0) circle (20bp);
    \node {\includegraphics[height=40bp,width=40bp,keepaspectratio]{#1}};
  \end{tikzpicture}%
}

% ——————————————————————————————————————————————————————————————————
% Весь блок автора
% ——————————————————————————————————————————————————————————————————
\newcommand{\authorfont}{\fontfamily{ftx}\selectfont}
\newcommand{\articleauthor}[3]{%
  % Три аргумента: фото, имя и ссылка на автора \href{url}{Имя}, дата
  \noindent
  \begin{minipage}[c]{40bp}
    \circlephoto{#1}
  \end{minipage}%
  \hspace{1em}%
  \begin{minipage}[c]{9cm}
    \authorfont\textbf{#2}\\
    {\small\authorfont\color{nnggray} #3}
  \end{minipage}
  \vspace{1em}\par
}

% ——————————————————————————————————————————————————————————————————
% Рамки через окружение
% ——————————————————————————————————————————————————————————————————
\renewenvironment{leftbar}{%
  \def\FrameCommand{{\color{gray!50}\vrule width 1pt}\hspace{10pt}}%
  \MakeFramed{\advance\hsize-\width \FrameRestore}%
  \noindent
}{%
  \endMakeFramed
}


% ——————————————————————————————————————————————————————————————————
% Оглавление без «Содержания»
% ——————————————————————————————————————————————————————————————————
\makeatletter
\renewcommand{\tableofcontents}{\@starttoc{toc}}
\makeatother